\section{引言}
中文显示测试：实验、测量、误差、频率、偏振、纵模、横模。五份文档使用相同内容，以便比较版面。文献引用使用自动编号\cite{lamport}。

\section{原理}
简谐函数的示例为
\begin{equation}
  y(t)=A\cos(\omega t+\varphi),\qquad f=\frac{\omega}{2\pi}.
  \label{eq:example}
\end{equation}
式\eqref{eq:example}中，$A$为振幅（与$y$同量纲），$\omega$为角频率（\si{\radian\per\second}），$\varphi$为初相位（\si{\radian}），$f$为频率（\si{\hertz}）。符号测试：$\alpha,\beta,\Delta,\mu,\Omega,\pm,\times$。

\section{实验}
本节仅介绍版式检查方法：查看中文、公式、表格和示意图，确认文字清楚、编号正确且没有裁切。本示例没有进行物理测量。

\section{结果与分析讨论}
表\ref{tab:symbols}展示符号和单位，仅作为排版示例。
\begin{table}[H]
  \centering\small
  \caption{物理量符号与单位示例，无实测数据。}
  \label{tab:symbols}
  \begin{tabularx}{\linewidth}{Xcc}
    \toprule
    物理量 & 符号 & 单位\\
    \midrule
    长度 & $L$ & \si{\metre}\\
    时间 & $t$ & \si{\second}\\
    频率 & $f$ & \si{\hertz}\\
    \bottomrule
  \end{tabularx}
\end{table}
Fig.~\ref{fig:example}展示解析余弦函数。相位增加$2\pi~\si{\radian}$时曲线完成一个周期，归一化振幅在$-1$与$1$之间变化。这是函数性质的示意，不是测量结果。
\begin{figure}[H]
  \centering
  % 单栏时控制示意图高度；双栏时保留足够宽度与可读字号。
  \newlength{\MVPDiagramWidth}
  \ifdim\linewidth>11cm\setlength{\MVPDiagramWidth}{0.65\linewidth}%
  \else\setlength{\MVPDiagramWidth}{0.9\linewidth}\fi
  \resizebox{\MVPDiagramWidth}{!}{%
  \begin{tikzpicture}[x=0.8cm,y=0.8cm,font=\small]
    \draw[->] (-0.1,0)--(7.2,0) node[right]{$\theta\,/\,\mathrm{rad}$};
    \draw[->] (0,-1.3)--(0,1.55) node[above]{$y/A$};
    \draw[thick,ReportAccent,domain=0:6.28,samples=80] plot (\x,{cos(\x r)});
    \node[anchor=south] at (3.4,1.08) {解析曲线：$y/A=\cos\theta$};
    \node[below left] at (0,0) {$0$};
    \node[below] at (6.28,0) {$2\pi$};
  \end{tikzpicture}}
  \caption{归一化余弦函数示意，横轴为相位（rad），纵轴为相对振幅（无量纲）；非实测数据。\newline\small Fig.~\thefigure\ Analytical cosine curve: phase in rad and dimensionless relative amplitude; no measured data.}
  \label{fig:example}
\end{figure}
排版检查不用于推断物理测量精度。

\section{结论和建议（总结）}
本示例用于选择样式。建议比较中文、公式、章节层级和图表尺寸。是否乱码应以实际查看效果为准。

\renewcommand{\refname}{六、参考文献}
\begin{thebibliography}{9}
  \bibitem{lamport} Lamport L. LaTeX: A Document Preparation System[M]. 2nd ed. Reading, MA: Addison-Wesley, 1994.
\end{thebibliography}
